\section{Terminology}
\label{app:terminology}

The following definitions specify how potentially overloaded terms are used throughout this paper.

\begin{description}[leftmargin=0pt,labelindent=0pt,itemsep=2pt,topsep=3pt]
  \item[Autoregressive language model (AR LM).]
  A language model that generates an output by predicting each next token from the input and the already generated prefix. Standard decoding advances left to right and treats earlier emitted tokens as fixed context.

  \item[Diffusion language model (dLLM).]
  A language model that generates through iterative denoising or refinement of a partially unknown or corrupted sequence. Depending on the architecture, several positions may be updated in one step and uncertain earlier positions may remain revisable.

  \item[Denoising step / diffusion update.]
  One iterative update of the diffusion state. A diffusion update can modify multiple positions or cognitive cells and is therefore distinct from generating one token in an AR model.

  \item[Reopening / remasking.]
  Increasing the editability of a previously predicted region so that later computation or newly arrived evidence can revise it. Remasking is one concrete token-level mechanism for reopening.

  \item[Tool call.]
  A discrete executable request to an external source, normally containing a tool or source identity and sufficiently bound arguments such as a query or structured JSON fields.

  \item[Observation.]
  Information returned by an external source or runtime event and made available to the model.

  \item[Information need.]
  A model state indicating that external information would be useful. In CID, such a need may become available to the runtime before it has converged into a complete serialized tool call.

  \item[Perceptual binding.]
  A persistent runtime relation connecting an information need to an external source, its current arguments, refresh policy, cached state, and affected cognitive regions.

  \item[External refresh.]
  Re-executing or re-reading an external source to obtain a potentially newer value. External refresh incurs external I/O or source computation.

  \item[Cognitive refresh / re-projection.]
  Recomputing how an already available external value should influence the current cognitive state, without repeating external I/O.

  \item[Perceptual projection.]
  A context-dependent model representation of an external value, constructed with respect to the current thought and display states and optionally linked to its provenance.

  \item[Fact channel.]
  Externally controlled information that the model may read but may not overwrite. Its contents can still change when the user, environment, runtime, or source updates them.

  \item[Thought channel.]
  The revisable internal cognitive state containing hypotheses, plans, information needs, percept interpretations, uncertainty, constraints, and intermediate conclusions.

  \item[Display channel.]
  The revisable token canvas that progressively converges to the user-visible response.

  \item[Typed Cognitive Tensor (TCT).]
  CID's structured continuous representation of the thought channel. Cognitive cells combine semantic vectors with typed roles, symbolic anchors, source links, uncertainty, local diffusion levels, and lifecycle state.

  \item[Local diffusion level.]
  A per-region value controlling how revisable that region currently is. Higher diffusion permits stronger revision; lower diffusion represents a more stable region.

  \item[Persistent perception.]
  The maintenance of an information source's cognitive influence while the corresponding need remains active. Persistence may reuse a static cached value or refresh a changing source according to policy.

  \item[Static and dynamic sources.]
  A static source is treated as unchanged while a binding is active, allowing repeated cognitive refresh from a cached value. A dynamic source can change over time and may require polling, version checks, subscriptions, or streaming updates.

  \item[Binding lifecycle.]
  The runtime state of a perceptual binding as it moves through stages such as \textsc{candidate}, \textsc{active}, \textsc{waiting}, \textsc{available}, \textsc{refreshing}, and \textsc{retired}.
\end{description}

\section{Reference Runtime Interface}
\label{app:runtime-interface}

A read-only source registered with a CID runtime may expose the following conceptual schema:

\begin{verbatim}
source:
  name: document_search
  description: >-
    retrieve passages from a corpus
  input_schema:
    query: text
    scope: optional identifier
    top_k: optional integer
  properties:
    read_only: true
    cacheable: true
    streamable: true
    cancellable: true
    versioned: true
  default_refresh:
    mode: on_version_change
    max_age_ms: 1000
\end{verbatim}

The runtime should not require that every source support every property. A calculator may be cacheable but not streamable; a clock is dynamic but inexpensive; a file may expose a version identifier; a search endpoint may return incremental results without a stable version.

\section{Binding State Machine}
\label{app:binding-state}

A perceptual binding can use the following lifecycle:
\begin{align}
  \textsc{candidate} &\rightarrow \textsc{active}
  \rightarrow \textsc{waiting} \notag\\
  &\rightarrow \textsc{available}
  \rightarrow \textsc{refreshing} \notag\\
  &\rightarrow \textsc{retired}.
\end{align}
A binding can move from \textsc{available} directly back to \textsc{active} when the value is unchanged and only cognitive re-projection is required. It moves to \textsc{refreshing} when the need remains active and the source freshness policy requires external work. Quiescence is a trajectory state rather than a binding state: active bindings remain intact while model updates pause, and the next observation resumes refinement without reconstructing the interaction context.

A candidate binding should record:
\begin{itemize}[leftmargin=*]
  \item source probability and need confidence;
  \item partially bound selector or arguments;
  \item links to originating cognitive cells;
  \item current cache key and source version;
  \item external and cognitive refresh policies;
  \item target fact, thought, and display regions.
\end{itemize}

\section{Conceptual Execution Traces}
\label{app:traces}

\subsection{Static Copying}

A user requests a response that must reproduce a value from a document exactly. The model forms a persistent need for the value and binds to a file reader. After the file is read once, the runtime caches the exact result. During subsequent denoising steps, the unchanged result is re-projected into the thought tensor, preserving salience while the display is reorganized. No repeated disk I/O is necessary, but the percept remains active until the exact value has been anchored in the display.

\subsection{Dynamic Display Length}

The current realized display length---for example, the inferred EOS position or populated span within the fixed canvas---is an externally computed value that the model may read but cannot rewrite. A binding to the length monitor remains active during generation. At each update, the runtime supplies a refreshed value, which changes the thought tensor's plan and density decisions. The display can compress earlier text because the value conditions both current cognition and already generated output regions.

\subsection{Streaming Retrieval}

A retrieval source first returns titles, then snippets, then complete passages. The binding remains active throughout the stream. Early chunks narrow hypotheses and refine the query; later chunks reopen unsupported conclusions and supply exact evidence. Unrelated display regions remain stable under local diffusion clocks.

\section{Suggested Training Data Construction}
\label{app:data}

A practical initial dataset can be synthesized from ordinary question--answer pairs and documents:
\begin{enumerate}[leftmargin=*]
  \item identify answer spans that depend on external evidence and hide those spans from the model-visible prompt;
  \item create source descriptors, binding targets, and delayed or incremental arrival schedules;
  \item construct pre-arrival states with unresolved cells and post-arrival revisions grounded in the source;
  \item add static-copy cases that preserve one value across many steps and dynamic cases in which the source changes before display convergence.
\end{enumerate}

Teacher-generated textual reasoning may be encoded into initial cognitive cells, but training should randomize cell ordering and event schedules to avoid reducing the thought tensor to a disguised text chain.

\section{Experimental Reproducibility}
\label{app:reproducibility}

\subsection{Model, Code, and Checkpoint Identity}

RQ1 uses \texttt{GSAI-ML/iLLaDA-8B-Base} at revision
\texttt{a1b5b5f8a31a3854a46205ee584178c04b45ec9a}. The StageA-v0 training code is pinned to commit
\texttt{a1396c79279b2aa024af1f97dad9130472c2dc25}. The frozen 8B backbone is reused for token
embeddings, the bidirectional decoder, and the native LM head; CID contributes trainable channel
projections, external-perception fusion, and prediction heads. The final StageA-v0 checkpoint contains
53 trainable tensors totaling 453,841,081 parameters. Its recorded state is 282,536 optimizer steps,
3,400,565 processed training transitions, and three completed epochs. The compact inference checkpoint
has SHA-256
\texttt{d9e33c65b4316b9c8acea93a53ebdf90}\allowbreak\texttt{afcf8d419d5c973c7d84dbc7ddb3b9c5}.
The checkpoint and all RQ1 raw artifacts are released at
\url{https://huggingface.co/fwerkor/CID-RQ1-StageA-v0}.

\subsection{CID-Dataset Training Corpus}

StageA-v0 is trained on the released \texttt{fwerkor/CID-Dataset} mixture. The materialized JSONL is
8,379,022,223 bytes and contains 192,297 semantic tasks expanded into 421,798 trajectories. There are
2,589,664 adjacent snapshot transitions plus 421,798 virtual empty-state bootstrap transitions, for
3,011,462 trainable transitions. Maximum trajectory length is 44 snapshots and the maximum required
TCT capacity is 128. The materialized file SHA-256 recorded by the release manifest is
\texttt{bec927029445815f085f5597ac1e4519}\allowbreak\texttt{6e8c64995793e218bc7f73df52eb8492}.
Semantic-task balancing is applied before component weights, and the global transition-weight mean is
normalized to one. Of the 192,297 semantic tasks, 134,750 (70.07\%) require tools, 43,933 (22.85\%)
require no tool, and 13,614 (7.08\%) expose tools that are unnecessary. The latter two groups provide
negative supervision against indiscriminate tool use.

\begin{table*}[t]
  \centering
  \caption{CID-Dataset components used by StageA-v0. ``Tasks'' are semantic tasks before schedule expansion; ``Traj.'' are materialized trajectories; ``Train trans.'' includes each component's bootstrap transitions. Weight is the semantic component multiplier applied after per-task balancing.}
  \label{tab:dataset-components}
  \scriptsize
  \begin{tabular}{lrrrr}
    \toprule
    Component & Tasks & Traj. & Train trans. & Weight \\
    \midrule
    public-base & 8,512 & 34,048 & 75,596 & 2.0 \\
    public-interaction & 7,590 & 30,360 & 430,834 & 2.5 \\
    mechanism & 10,000 & 40,000 & 214,287 & 1.0 \\
    computational & 12,000 & 24,000 & 135,339 & 1.0 \\
    symbolic & 15,600 & 31,200 & 153,104 & 1.0 \\
    correction & 10,000 & 20,000 & 140,220 & 1.0 \\
    complex-logic & 12,000 & 24,000 & 144,000 & 1.0 \\
    self-identity & 720 & 1,440 & 2,880 & 1.0 \\
    multilingual & 1,200 & 2,400 & 21,725 & 1.0 \\
    composed & 12,000 & 24,000 & 261,058 & 1.0 \\
    tool-restraint & 6,500 & 6,500 & 6,500 & 1.5 \\
    compositional-longtail & 20,000 & 40,000 & 220,000 & 1.0 \\
    long-horizon & 12,000 & 24,000 & 471,626 & 1.0 \\
    deep-tool-restraint & 4,000 & 4,000 & 26,000 & 1.5 \\
    natural-interaction & 8,675 & 17,350 & 182,059 & 2.0 \\
    natural-public/nq-open & 30,000 & 60,000 & 330,121 & 1.0 \\
    natural-public/oasst1 & 4,500 & 4,500 & 9,000 & 1.5 \\
    natural-public/multidoc2dial & 15,000 & 30,000 & 165,171 & 1.5 \\
    natural-public/qasper & 2,000 & 4,000 & 21,942 & 2.0 \\
    \midrule
    Total & 192,297 & 421,798 & 3,011,462 & -- \\
    \bottomrule
  \end{tabular}
\end{table*}

The natural-public subset contributes 30,000 NQ-Open train tasks (CC-BY-SA-3.0), 4,500 OASST1 train
tasks (Apache-2.0), 15,000 MultiDoc2Dial train tasks (CC-BY-3.0), and 2,000 QASPER train tasks
(CC-BY-4.0). Public-source rows retain upstream repository, revision, split, license, and row-key
provenance. Teacher-generated components contain concise semantic state summaries rather than private
chain-of-thought transcripts. Physical slot placement and event timing are schedule variables rather
than teacher outputs.

\subsection{StageA-v0 Optimization}

\begin{table}[h]
  \centering
  \caption{StageA-v0 optimization and model limits. The effective global batch is held at 32 when the run is resumed with fewer GPUs.}
  \label{tab:stagea-hparams}
  \scriptsize
  \begin{tabular}{ll}
    \toprule
    Setting & Value \\
    \midrule
    Optimizer & AdamW \\
    Peak learning rate & $1\times10^{-4}$ \\
    Weight decay & 0.01 \\
    Precision & bfloat16 \\
    Micro-batch / GPU & 1 \\
    Effective global batch & 32 \\
    Gradient accumulation & 8 (4 GPUs), 16 (2 GPUs) \\
    Epochs & 3 \\
    Warmup steps & 2,000 \\
    LR decay steps & 285,000 \\
    Minimum LR ratio & 0.1 \\
    Gradient-norm clip & 1.0 \\
    Diffusion $t$ range & $[0.05,1.0]$ \\
    Rollout horizon & 3 \\
    Teacher-forcing epochs & 1 \\
    Rollout-ramp epochs & 2 \\
    Maximum TCT slots & 128 \\
    Maximum display tokens & 1,024 \\
    Runtime allocation threshold & 0.8 \\
    Runtime need threshold & 0.6 \\
    \bottomrule
  \end{tabular}
\end{table}

StageA-v0 was run with data-parallel RTX A6000 48\,GB GPUs and resumed across world sizes while
preserving effective batch 32. The final phase used two GPUs with gradient accumulation 16; earlier
four-GPU phases used accumulation 8. The backbone parameters have \texttt{requires\_grad=False}, but
the forward and input-gradient path still traverses the full decoder because CID projections precede
it. Thus Stage A is materially cheaper in optimizer state and parameter-gradient computation than
Stage B, but it is not equivalent to training only a small standalone 454M network.

The loss weights in StageA-v0 are: thought 1.0, display 1.0, convergence 0.2, allocation 0.3,
roles 0.2, uncertainty 0.1, local-noise delta 0.1, lifecycle 0.1, information-need intent 0.5,
source selection 0.5, argument presence 0.1, argument grounding 0.2, revision 0.2, refresh 0.2,
anchor presence/kind/grounding 0.1/0.1/0.2, and link presence/relation/target-kind/grounding
0.1/0.1/0.1/0.2. Allocation uses masked binary cross-entropy over currently empty physical slots;
this exact formulation is the RQ1 defect analyzed in \cref{sec:rq1}.

\subsection{RQ1 Held-Out Protocol}

The RQ1 NQ set contains 32 rows sampled with seed 20260824 from
\texttt{google-research-datasets/nq\_open}, revision
\texttt{5dd9790a83002ad084ddeb7c420dc716852c6f28}, file
\texttt{nq\_open/validation-00000-of-00001.parquet}. Training uses the upstream train split only.
Each held-out trajectory has one \texttt{knowledge\_search} evidence event with a randomized one-to-three
step delay and an eight-slot teacher schedule. The default-runtime control uses 128 TCT slots, four
denoising iterations per runtime update, and a 16-step cap. Diagnostic runs use the same weights,
source replay, argument catalog, denoising budget, and task set. Runs are executed in bfloat16 on RTX
3090 24\,GB GPUs, one model replica per probe job. Raw outputs, the exact held-out JSONL, registry,
validation scripts, aggregate summary, and SHA-256 manifest are included in the public RQ1 artifact.

We additionally use a ten-task compositional generalization probe with zero exact logic-spec overlap
with training. This probe spans boolean, graph, spatial, causal, ordering, quantifier, candidate,
numeric, policy, and repair tasks and is used only as a checkpoint-health control, not as a general
benchmark claim.

\subsection{RQ3--RQ6 Protocol}

Task families cover static reference/copying, dynamic state tracking, delayed retrieval, streaming
evidence, and competing sources. Baselines include blocking ReAct; an event-driven asynchronous AR
agent \citep{ginart2024asynchronous,hooper2026realtime}; a dLLM with conventional explicit-call
regions; P-ReAct/DLLM-Searcher \citep{zhao2026dllmsearcher}; one-time CID percept injection;
persistent static re-projection; full CID with static/dynamic refresh; and oracle bindings as an
upper bound.

We report exact/task-specific correctness, evidence support, copying fidelity, stale-value rate,
wall-clock latency, tool-wait overlap, binding precision/recall, cache hits, external refreshes,
duplicate-I/O suppression, repeated cognitive projection, and revision locality. Intent lead time is
$\Delta_{\mathrm{lead}}=t_{\mathrm{explicit\ call}}-t_{\mathrm{binding}}$, and assimilation lag is
$\Delta_{\mathrm{assim}}=t_{\mathrm{correct\ state}}-t_{\mathrm{arrival}}$. Quality comparisons use
a logical event clock so accelerator speed cannot change pre-arrival refinement; matched-hardware
wall-clock measurements are reported separately. Ablations remove role typing, symbolic anchors,
source links, per-cell noise, persistent bindings, static re-projection, dynamic refresh,
display-to-thought feedback, and arrival-time randomization. Where meaningful, the three RQ1
contract defects are retained as explicit negative ablations rather than disappearing after repair.

\subsection{RQ2 and Stage B Reporting Rule}

RQ2 retrains Stage A from the pinned base model after the contract fixes described in
\cref{sec:rq2}; it does not initialize from StageA-v0. Stage B then unfreezes the backbone and uses
FSDP/AdamW. To avoid reporting stale pilot settings as final hyperparameters, the RQ2--RQ6 table is
populated only from the frozen retrained checkpoint metadata after the run completes. The earlier
pre-fix Stage-B smoke/stress runs are engineering tests and are excluded from experimental results.
This rule is important because the RQ1 findings can change the optimal capacity, display, and
allocation settings.

\section{Open Technical Questions}
\label{app:questions}

Several design decisions remain open:
\begin{itemize}[leftmargin=*]
  \item which iterative process should update the thought tensor, how many cells it needs, and whether cells should be created or retired dynamically;
  \item whether display and thought diffusion should share a backbone and how source descriptors generalize to unseen tools;
  \item how to estimate support and conflict for local reopening while preserving privacy during debugging;
  \item how to tune refinement-epoch and wall-clock budgets across deployment settings.
\end{itemize}
These questions remain open design variables within the CID framework.
